\documentclass[11pt,a4paper]{article}
\usepackage{lmodern}
\usepackage[utf8]{inputenc}
\usepackage[T1]{fontenc}
\usepackage{amsmath}
\usepackage{amsfonts}
\usepackage{amssymb}
\usepackage{graphicx}
\usepackage[french]{babel}
\usepackage[notitlepage,fancysections,pagenumber]{sauveterre}

\usepackage{enumitem}

\usepackage{indentfirst}
\usepackage{lipsum}

\usepackage{cases}
\usepackage{multicol}

%formatting titles
\usepackage{titlesec}
\titleformat*{\section}{\centering\Large\bfseries}
\titleformat*{\subsection}{\centering\large\bfseries}

%formatting date
\usepackage[ddmmyyyy]{datetime}
\renewcommand{\dateseparator}{/}

%\usepackage{geometry}
%\geometry{
% a4paper,
% total={170mm,257mm},
% left=20mm,
% top=20mm,
% }

\usepackage{pgf,tikz}
\usepackage{mathrsfs}
\usetikzlibrary{arrows}

\usepackage[export]{adjustbox}

\renewcommand{\notitlepageupperline}{
   \noindent%
   \begin{minipage}{\textwidth}
      \centering
      \begin{minipage}{0.33\textwidth}
         \raisebox{-0.5\height}
         {\color{bleu-sauveterre}%
         \includegraphics[height=0.8\hauteurlogovert]
         {\polylogovert}
         }
      \end{minipage}%
      \begin{minipage}{0.33\textwidth}
         \color{bleu-sauveterre}%
         \centering\polydatesave
         \polyauthorsave
         \hspace{3em}
      \end{minipage}%
      \begin{minipage}{0.33\textwidth}
         \color{bleu-sauveterre}%
         \flushright\raisebox{-0.5\height}
         {\includegraphics[height=0.8\hauteurlogovert]{logo-federal}
         }
      \end{minipage}
   \end{minipage}
}

%\usepackage{fancyhdr}
%\usepackage{lastpage}
%\pagestyle{fancy}
%\fancyhf{}
%%headers and footers
%\lhead{Centre Équestre de Sauveterre}
%\chead{Règlement intérieur}
%\rhead{Version du \today}
%\rfoot{Page \thepage \ sur \pageref{LastPage}}

%\fancypagestyle{firstpage}{%
%  \fancyhf{}
%  \renewcommand\headrulewidth{0pt}
%  \fancyfoot[R]{Page \thepage \ sur \pageref{LastPage}}
%}
%
%\title{\vspace{-50pt}
%\includegraphics[width=0.4\textwidth , left]{logo2-v2.eps}~
%\\[8pt]
%{\huge \textsc{Règlement intérieur}}
%}

\title[Conditions Générales applicables aux Balades \\ version du \today]{Conditions Générales applicables aux Balades}
\subtitle{Annexe I au \textit{Règlement Intérieur} (\no 24--09/CES/DIR)}
\author{}
\date{}

%enumerate environment formatting
\renewcommand{\labelenumi}{\textbf{\arabic{enumi}.}}
\renewcommand{\labelenumii}{\textbf{\arabic{enumi}.\arabic{enumii}.}}
\renewcommand{\labelenumiii}{\textbf{\arabic{enumi}.\arabic{enumii}.\arabic{enumiii}.}}

%change sec and subsec numbering
\renewcommand{\thesection}{Section \Roman{section}.}
\renewcommand{\thesubsection}{Article \arabic{subsection} :}

\renewcommand{\labelenumi}{\arabic{enumi}.}


\renewcommand*\contentsname{Plan}

\renewcommand{\emph}[1]{\textit{\textbf{#1}}}


\usepackage{framed}

%\setlength{\parindent}{0pt}%remove indentation

\begin{document}

   \maketitle

%\thispagestyle{firstpage}
   \vspace{-30pt}

   \subsection*{Balades à poney}
      Les tours de poneys en main se font sous l'entière responsabilité des parents, de la famille, ou de leur délégué.

      L'accompagnant familial doit disposer de toutes ses facultés physiques pour assurer la sécurité de l'enfant à poney lors de la prise en main.
      Tout handicap reconnu (articulations foulées ou cassées\dots) ou situation ponctuelle (femme enceinte\dots) interdit la prise en main pour des raisons évidentes de sécurité.

      Il est demandé aux familles de prévoir \emph{1 adulte par poney}.
      Un départ en balade peut être refusé si le nombre d'adultes est insuffisant pour tenir les poneys.

      Les balades en main se font \emph{uniquement au pas}, à raison d'un enfant seulement par poney, ce dernier devant obligatoirement porter une bombe.
      Le non-respect de ces conditions peut entrainer une exclusion du Centre Équestre comme prévu dans l'article 8 du \textit{Règlement Intérieur}.

   \subsection*{Balades à cheval}
      Les balades à cheval se font seulement sur réservation (téléphone, mail, sur place\dots).
      Il peut être demandé, pour des groupes d'effectif supérieur à 2 personnes, un règlement de montant au moins égal à 1 place (au tarif en vigueur à la réservation).
      \\

      Toute balade à cheval est encadrée par le nombre d'enseignants défini par la direction, en accord avec les règlements et autorisations afférents aux diplômes et assurances.
      Ces derniers sont alors responsables de leur groupe, et sont habilités à prendre toute décision nécessaire au maintien de la sécurité aussi bien des cavaliers que du groupe dans son ensemble.
      \\

      On distingue, en balade, deux groupes distincts : les \og débutants \fg{} (toute personne non expérimentée montant à cheval pour la première fois ou seulement occasionnellement, sans niveau fédéral) et les \og confirmés \fg{} qui sont les cavaliers possédant un niveau équestre (du point de vue fédéral) et étant en capacités de trotter et galoper en toute aisance en extérieur\footnote{Minimum Galop 3 selon les programmes fédéraux.}.
      Tout cavalier ne répondant pas à ces critères est \textit{de facto} considéré comme débutant.

      Le groupe constitué des cavaliers débutants fera la balade \emph{au pas uniquement}.
      Au contraire une balade de niveau confirmé comportera des phases de trot et de galop, \emph{de longueur et durée décidées par l'enseignant qui est le seul à avoir connaissance des capacités et de l'état de sa cavalerie} et du niveau des cavaliers qu'il encadre.

      Les enseignants ayant pour mission première de garantir la sécurité du groupe (\textit{cf. supra}), ils ont la possibilité de refuser des cavaliers dans le groupe confirmé s'ils estiment que trotter ou galoper présente un risque pour leur sécurité. \emph{La responsabilité des enseignants ne saurait être engagée en cas de chute d'un cavalier n'ayant pas respecté cette décision, ou toute autre consigne de sécurité donnée avant ou pendant la balade}.

   \subsection*{Tenue exigée et aptitude physique}
      Aucune tenue spécifique n'est exigée, cependant il est conseillé aux cavaliers, pour leur sécurité et leur confort, de porter un pantalon long et des chaussures fermées.
      Malgré tout, l'enseignant responsable du groupe se réserve le droit, en accord avec la direction, de refuser un cavalier portant une tenue non adaptée à la pratique de l'équitation en extérieur et ne permettant pas d'assurer sa sécurité ou celle du groupe (ex : tongs, sac à dos, écharpe, déguisement\dots).\\

      Quiconque monte en balade, quel que soit son niveau, est pleinement conscient qu'il s'agit d'une pratique sportive d'extérieur.
      Aussi, il est de la responsabilité de chacun de s'assurer qu'il est en pleine capacité de pratiquer l'équitation, même au pas et en balade.

      En particulier, \emph{toute condition médicale ou physique incompatible avec l'activité proposée est supposée connue du pratiquant}, qui monte alors en pleine conscience des risques éventuels du fait de sa condition.
      En cas de doute sur la pleine capacité à pratiquer l'équitation en extérieur, les enseignants peuvent refuser le départ d'un cavalier, pour des raisons de sécurité individuelle et collective.
      Même si nous faisons tout pour le réduire, \emph{le risque de chute existe}.

   \subsection*{Handicap}
      Nous défendons l'équitation pour tous, et nous engageons à proposer à chacun une activité adaptée à ses capacités.
      À ce titre, il découle du bon sens que \emph{tout handicap ou situation nécessitant un accompagnement personnalisé doit être signalé au moment de la réservation}, afin que nous puissions nous organiser au mieux pour vous accueillir et adapter la cavalerie.

      Pour des raisons de sécurité et de bien-être animal, la direction se réserve le droit de refuser un cavalier n'ayant pas rempli la condition précitée en amont de la prestation.

   \subsection*{Annulation et remboursement de la prestation}
      Toute balade à cheval réservée peut être annulée jusqu'à \emph{12 heures} avant l'heure de rendez-vous.
      Passé ce délai, les sommes avancées (paiement intégral ou partiel) ne pourront être remboursées, hors cas de force majeure (urgence médicale ou familiale, sur présentation d'un justificatif).

      Concernant les balades à poney, la réservation n'étant pas obligatoire, toute balade payée à l'avance pourra être remboursée si elle est annulée, sur présentation du justificatif donné lors du règlement (ticket, facture\dots).
      \\

      Un faux-départ n'est pas remboursé.
      Aussi, toute balade commencée est due.
      Une balade est considérée comme \og commencée \fg{} à partir du moment où le cavalier est à cheval, et le groupe à l'extérieur du manège.
      De même, une balade à poney n'est plus remboursable à partir du moment où l'enfant à poney entre dans la forêt.

      Aucun remboursement n'est possible pour un cavalier auquel aurait été refusé le galop.

      Les cartes cadeau, quelles que soient leur valeur et leur validité, ne pourront pas être remboursées ni échangées.
      De plus, leur validité ne peut être repoussée hors cas de force majeure.

      Dans les autres cas, le remboursement est décidé par la direction.


      \vspace{3cm}

      \hspace{.55\textwidth}
      \begin{minipage}{.4\textwidth}
         \begin{center}
            Stéphanie \textsc{Marais},\\
            Dirigeant le Centre Équestre de Sauveterre\\
            \textbf{\texttt{Original Signé}}
         \end{center}
      \end{minipage}

\end{document}
